\subsection{CT15 --- Rank forcing}\label{ct:15}\label{rank-forcing-charge-only-independent-coordinates}
\ctsignature{\(P_{\ast\to15}\to \ctcert{C4};\ \ctint{P_{15\to3},P_{15\to7},P_{15\to16},P_{15\to4}};\ \ctrec{\ctnone}.\)}

CT15 is a target-relative rank firewall.  Its entry contract fixes the test
family, rank map, target datum, capacity datum, branch state, and C4 claim.
The executable framework separately states that the rank is target-relative,
that a drop yields structural dependence, that full rank yields independent
coordinates, that the ledger is certified, and that rank demand exceeds
capacity.

The rank certificate has one successor.  The rank-drop decision then produces
either a \texttt{DependenceState} or a \texttt{FullRankState}.  Structural
dependence enters a dedicated routing node and constructs an exact CT3, CT7,
or CT16 input.  Full rank enters a dedicated ledger certificate; the comparison
node then closes by C4 or constructs an exact CT4 input.  No ledger can be
constructed on the dependence branch, so incidental rank cannot consume
capacity.

\begin{itemize}
\tightlist
\item \textbf{S-Def} is the target-relative rank certificate.
\item \textbf{S-Dich} is the rank-drop decision.
\item \textbf{S-Comp} is the independent-coordinate ledger and comparison.
\item \textbf{S-Rout} and \textbf{S-Trig} are exact CT3, CT4, CT7, and CT16 inputs.
\end{itemize}

\begin{figure}[p]
\centering
\resizebox{0.98\textwidth}{!}{%
\begin{tikzpicture}[x=1cm,y=1cm,every node/.style={font=\scriptsize}]
\node[bcwide] (entry) at (0,0) {{\tiny\ttfamily CT15.entry}\\\textbf{Validated rank datum}};
\node[bcdec,bclemma] (scope) at (0,-2.1) {{\tiny\ttfamily CT15.decide.scope}\\\textbf{Finite target-relative rank?}};
\node[bcscopefail] (scopeout) at (-7.4,-2.1) {{\tiny\ttfamily CT15.terminal.scope}\\\textbf{Scope exit}};
\node[bcwide,bclemma] (rank) at (0,-4.4) {{\tiny\ttfamily CT15.certify.rank}\\\textbf{Target meaning of rank}};
\node[bcdec,bclemma] (drop) at (0,-6.8) {{\tiny\ttfamily CT15.decide.rankDrop}\\\textbf{Dependence or full rank?}};
\node[bcroutechoice,bctheorem] (route) at (-5.6,-9.3) {{\tiny\ttfamily CT15.decide.\allowbreak dependenceRouting}\\\textbf{Classify dependence}};
\node[bcroute,bcrouteneutral] (ct3) at (-10.0,-12.0) {{\tiny\ttfamily CT15.terminal.ct3}\\\textbf{CT 3 exit}};
\node[bcroute,bcrouteneutral] (ct7) at (-5.6,-12.0) {{\tiny\ttfamily CT15.terminal.ct7}\\\textbf{CT 7 exit}};
\node[bcroute,bcrouteneutral] (ct16) at (-1.2,-12.0) {{\tiny\ttfamily CT15.terminal.ct16}\\\textbf{CT 16 exit}};
\node[bcwide,bclemma] (ledger) at (5.6,-9.3) {{\tiny\ttfamily CT15.certify.ledger}\\\textbf{Independent-coordinate ledger}};
\node[bcdec,bclemma] (compare) at (5.6,-12.0) {{\tiny\ttfamily CT15.decide.comparison}\\\textbf{Demand exceeds capacity?}};
\node[bcclosed] (c4) at (3.2,-14.6) {{\tiny\ttfamily CT15.terminal.c4}\\\textbf{C4}};
\node[bcroute,bcrouteneutral] (ct4) at (8.0,-14.6) {{\tiny\ttfamily CT15.terminal.ct4}\\\textbf{CT 4 exit}};
\draw[bcarr] (entry)--(scope); \draw[bcscopearr] (scope)--(scopeout);
\draw[bcarr] (scope)--(rank); \draw[bcarr] (rank)--(drop);
\draw[bcarr] (drop)--node[bclabel,above,sloped]{dependent}(route);
\draw[bchandoff] (route)--(ct3); \draw[bchandoff] (route)--(ct7); \draw[bchandoff] (route)--(ct16);
\draw[bcarr] (drop)--node[bclabel,above,sloped]{full rank}(ledger); \draw[bcarr] (ledger)--(compare);
\draw[bcarr] (compare)--(c4); \draw[bchandoff] (compare)--(ct4);
\end{tikzpicture}%
}
\caption[Closure-tactic 15: exact rank-forcing machine]{CT15 makes the dependence and full-rank branches disjoint at the type level; only full rank can reach ledger certification.}
\label{fig:ct15-rank-forcing}
\end{figure}
\FloatBarrier
